% A.62 candidate. Notes only until full independent audits.
\subsection{The ordinary-prime Heisenberg family by actual isogeny descent}

\begin{statement}{A.62}{Proved}{Both assertions for every ordinary-prime Heisenberg cover}
For every prime \(p\equiv1\pmod4\), let \(C_p\) be the smooth
proper normalization of \(z^p=\psi_p(x)\) over
\(E:y^2=x^3-x\), and let
\[
 Y_p=C_p\setminus q^{-1}(E[p])\xrightarrow{[p]q}U.
                                                        \tag{A.62.1}
\]
This actual degree-\(p^3\) Heisenberg cover has proper genus
\(1+p^2(p-1)/2\) and \(p^2\) cusps. Both assertions of
Task A.5 hold at every coefficient prime and for every
permitted cusp packet and open derived subgroup.

At a place above \(p\), its proper stable reduction has one
ordinary elliptic main component, \(p\) tails of the form
\(V^p-V=U^4\), and \(p(p-3)\) tails of the form
\(V^p-V=U^2\), over the algebraic residue field. Every
remaining component is rational and the dual graph is a tree.
The \(p^2\) marks form \(p\) rational clusters of \(p\)
marks each, disjoint from every positive-genus tail.

Together with A.61, this proves both all-prime assertions for
the actual exponent-\(p\) Heisenberg cover at every odd prime
\(p\). It makes no cofinality assertion about all finite covers.
\end{statement}

\begin{proof}
All special-fibre functions below are initially defined over
\(\F_p\); finite unramified extensions will separate the
points used in the actual charts. Since \(p\equiv1\pmod4\),
one has \(p\ge5\).

\paragraph{The separable isogeny and the actual reduced division polynomial.}
The coefficient of \(x^{p-1}\) in
\((x^3-x)^{(p-1)/2}\) is, up to sign,
\(\binom{(p-1)/2}{(p-1)/4}\), and is nonzero modulo \(p\).
Thus the good elliptic reduction is ordinary by
\cite[Theorem V.4.1(a), p.148]{fghSilvermanAEC}.
Its reduced \(p\)-torsion subgroup \(H\) has order \(p\),
and the inseparable degree of \([p]\) is \(p\).

There is an actual degree-\(p\) CM endomorphism
\(\alpha:E\to E\), over \(\Z_p\), with unit differential.
Indeed choose a Gaussian factor of \(p\) that is a unit
under the chosen embedding \(\Z[i]\to\Z_p\), and apply
that endomorphism to the smooth elliptic group scheme.
Duality sends the degree-one automorphism \([i]\)
to its inverse \([-i]\) and is additive. Hence the
dual of this Gaussian endomorphism is its conjugate;
their composite is \([p]\). The degree is therefore
\(p\), by \cite[Theorem III.6.2(a),(c),(d), p.83]{fghSilvermanAEC}.
Its kernel is finite etale over \(\Z_p\), since its
differential is a unit, and its special kernel is exactly
\(H\). Write \(\overline\alpha^*\omega=\mu\omega\),
where \(\omega=dx/(2y)\) and \(\mu\in\F_p^\times\).
The endomorphism commutes with the quarter-rotation
\(\rho(x,y)=(-x,iy)\).

The exact numerator and denominator in
\[
 x\circ[p]=\phi_p/\psi_p^2
                                                        \tag{A.62.2}
\]
are relatively prime after reduction; \(\phi_p\) is monic
of degree \(p^2\), by
\cite[Exercise III.3.7(a)--(d), pp.105--106]{fghSilvermanAEC}.
At every point of \(H\), the special map \([p]\) has
ramification index \(p\). At each nonzero such point,
the denominator therefore has a zero of order \(2p\).
At \(O\), the pole order of the quotient is \(2p\).
Comparing degrees in (A.62.2) gives
\[
 \overline\psi_p=\overline g^p,\qquad
 \deg\overline g=(p-1)/2,\qquad
 \operatorname{div}(\overline g)=D_H-p(O).
                                                        \tag{A.62.3}
\]
Here the coefficients have unique Frobenius roots in
\(\F_p\). The zeros of \(\overline g\) are the nonzero
points of \(H\), each simple. As in A.61, \(\psi_p\)
is even in \(x\); therefore so is \(\overline g\).
Choose an even integral polynomial lift \(g\) of the
same degree and set
\[
 h=\frac{\psi_p-g^p}{p}\in\Z_p[x],\qquad
 H_0=\frac{h}{g^p},\qquad
 \frac{\psi_p}{g^p}=1+pH_0.
                                                        \tag{A.62.4}
\]
The subscript distinguishes this rational function from
the reduced torsion group. A different integral lift of
\(\overline g\) gives the same \(\overline h\), since
\((g+pb)^p\equiv g^p\pmod{p^2}\).
The degree of \(h\) is \((p^2-1)/2\), with leading
coefficient one. In particular \(\overline H_0\) has
an exact pole of order \(p-1\) at \(O\).

\paragraph{Actual translation comparison modulo p squared.}
We prove that \(d\overline H_0\) is invariant under
every translation in \(H\). This comparison uses the
actual division function, not only its branch divisor
after reduction.

After a finite unramified extension \(R/\Z_p\), every
point of \(\ker\alpha\) is an \(R\)-section. Fix one,
\(T\), and write \(K=\operatorname{Frac}(R)\).
There is a function \(k_T\in K(E)^\times\) with
\(\operatorname{div}(k_T)=p(O)-p(-T)\), since \(pT=0\).
The characteristic-zero division divisor consequently gives
\[
 \frac{\psi_p(P+T)}{\psi_p(P)}=c_T k_T(P)^p,
                         \qquad c_T\in K^\times.
                                                        \tag{A.62.5}
\]
Let \(v_G\) be the Gauss valuation at the generic point
of the smooth elliptic special fibre. Its value group
on \(K(E)\) is \(\Z\), with uniformizer \(p\).
Scale \(k_T\) by a power of \(p\) so that \(v_G(k_T)=0\).
The two division functions in (A.62.5) are Gauss units,
by (A.62.3), so \(c_T\in R^\times\).

Choose \(\gamma_T\in R^\times\) with
\(\overline\gamma_T^p=\overline c_T\), possible
because the residue field is perfect. Reduction of
(A.62.5) and injectivity of Frobenius in its function
field show that
\[
 q_T(P)=\frac{g(P+T)}{g(P)k_T(P)\gamma_T}
                  \equiv1\pmod p
                       \quad\text{in the Gauss DVR}.
                                                        \tag{A.62.6}
\]
Thus \(q_T=1+pb_T\) with \(b_T\) in that DVR, and
\(q_T^p\equiv1\pmod{p^2}\). Also
\(c_T/\gamma_T^p=1+p\epsilon_T\), with
\(\epsilon_T\in R\). The exact equations now give
\[
 1+pH_0(P+T)
   =\frac{c_T}{\gamma_T^p q_T(P)^p}(1+pH_0(P)),
 \qquad
 \overline H_0(P+\overline T)
          =\overline H_0(P)+\overline\epsilon_T.
                                                        \tag{A.62.7}
\]
All divisibilities here take place over an unramified
field, so reduction modulo \(p^2\) has its literal
meaning. The possibly nonzero constant in the last
line is harmless; its differential is zero.
The residues of these sections run through all of
\(H\), proving the desired invariance.

\paragraph{Descent, pole order, and every critical zero.}
Let \(D\) be the invariant derivation
\(Dx=2y\), \(Dy=3x^2-1\), and put
\[
        F=D\overline H_0,
        \qquad U=D\overline g/\overline g.
                                                        \tag{A.62.8}
\]
The exact characteristic-zero identity (A.61.11), with
\(\psi_p=g^p+ph\), gives on division by \(p\) and
reduction
\[
 D^2\overline h
   =\overline g^{p-2}(D\overline g)^2
      -\overline g^{p-1}D^2\overline g,
                \qquad DF=-DU.                       \tag{A.62.9}
\]
Translation invariance from (A.62.7) proves that the
actual function \(F\) descends through
\(\overline\alpha\): write \(F=f\circ\overline\alpha\).

We also determine \(DU\) exactly. For \(T\in H\),
the divisor of \(\overline g(P+T)/\overline g(P)\)
is \(p(O)-p(-T)\). Its logarithmic differential has
no pole, because all these orders are divisible by
\(p\). It is therefore a constant multiple of
\(\omega\). Thus \(U(P+T)-U(P)\) is constant and
\(DU\) is \(H\)-invariant. It has at most double
poles, precisely over \(H\), so its descent belongs
to \(L(2O)=\langle1,x\rangle\) on the quotient elliptic
curve, identified with \(E\) by \(\overline\alpha\).
The quarter-rotation sends \(DU\) to \(-DU\),
because \(\overline g\) is invariant and \(\rho\)
scales \(D\) by \(-i\). It sends the quotient
\(x\)-coordinate to \(-x\). The constant term in
that two-dimensional space must consequently vanish.

At \(O\), the pole order of \(\overline g\) is
\(p-1\), so in an invariant local parameter to
the required first order \(U=s^{-1}+O(1)\) and
\(DU=-s^{-2}+O(s^{-1})\). Since
\(x\circ\overline\alpha=\mu^{-2}s^{-2}+O(s^{-1})\),
the leading coefficient determines
\[
       DU=-\mu^2x\circ\overline\alpha,
          \qquad Df=\mu x.                           \tag{A.62.10}
\]
Only first-order local expansions are used here;
no characteristic-\(p\) formal logarithm is assumed.

The only possible poles of \(F\) lie in \(H\).
Its pole at \(O\) has exact order \(p\), by the
pole order \(p-1\) of \(\overline H_0\), which is
not divisible by \(p\). The isogeny is etale;
therefore \(f\) has exactly one pole, at \(O\),
of exact order \(p\). The function \(F\), and
hence \(f\), is odd under \([-1]\).
In particular \(f(0,0)=0\). Equation (A.62.10)
has a zero of order two there, so \(f\) has a
zero of order three. Indeed \(3<p\), and the
first nonzero differentiated coefficient is of
degree two. At every other zero of \(f\), the
right side is nonzero, so the zero is simple.
Its zero divisor has total degree \(p\). Thus
there is one triple zero and \(p-3\) simple zeros.
Pullback through the etale degree-\(p\) isogeny gives
exactly \(p\) triple and \(p(p-3)\) simple zeros
of \(F\), all outside \(H\).

The differential \(d\overline H_0=F\omega\)
has these same zero orders. At a zero of order
\(k-1\), with \(k=4\) or \(2\), its primitive
\(\overline H_0-\overline H_0(P)\) has exact
order \(k\): every smaller degree is less than
\(p\), so differentiation detects it. In particular
there is no hidden lower degree-\(p\) coefficient
in this case. The two-torsion centres use \(y\)
as parameter; at all others \(x-x(P)\) is a parameter.

\paragraph{Exact main and all tail charts.}
Choose \(\pi^p=p\) and set \(z=g(x)+\pi u\).
The exact equation, divided by \(p\), has special
function field
\[
                         u^p=\overline h(x).
                                                        \tag{A.62.11}
\]
Its differential is nonzero because \(F\ne0\).
This is the degree-\(p\) radicial field
\(k(E)^{1/p}\), with ordinary elliptic normalization.
The chart is smooth at its special generic point,
by the differential and the elliptic equation.
It therefore supplies an actual elliptic component,
not just a formal equation.

Separate all critical centres over a finite
unramified extension of \(\Z_p\). At a non-two-torsion
centre lift its \(x\)-coordinate arbitrarily in that
unramified ring and lift \(y\) by Hensel's lemma;
its \(y\)-coordinate is a unit. At a two-torsion
centre use \((a,0)\), \(a=0,1,-1\), itself.
Every such \(g(a)\) is a unit, since the centres
avoid \(H\). Set
\[
 \begin{aligned}
 H_a&=h(a)/g(a)^p,&
 N_a(x)&=h(x)-H_a g(x)^p,\\
 F_a&=1+pH_a,& z_a^p&=F_a,\\
 \eta^{p-1}&=-p,& \rho^k&=\eta,
                \qquad z=g(x)z_a(1+\eta V).
 \end{aligned}                                         \tag{A.62.12}
\]
The exact binomial equation is
\[
 \begin{aligned}
 \mathcal B_\eta(V)
   &=\sum_{j=1}^{p-1}\frac{\binom pj}{p}\eta^{j-1}V^j-V^p,\\
 g(x)^pF_a\mathcal B_\eta(V)&=N_a(x)/\eta.
 \end{aligned}                                         \tag{A.62.13}
\]
At a non-two-torsion lift \((a,y_a)\), put
\(x=a+\rho U\), \(y=y_a+\rho W\). The exact
elliptic equation is
\[
 2y_aW+\rho W^2
          =(3a^2-1)U+3a\rho U^2+\rho^2U^3.
                                                        \tag{A.62.14}
\]
The first nonconstant Taylor degree of
\(\overline N_a\) is \(k\), with a nonzero
coefficient. Every smaller integral coefficient
is divisible by \(p\), because the lifts are
unramified. After division by \(\eta\), a lower
degree \(j<k\) has valuation at least
\[
                 1+\frac{j-k}{(p-1)k}>0;
                                                        \tag{A.62.15}
\]
a higher degree has valuation at least
\((j-k)/((p-1)k)>0\). Thus only degree \(k\)
survives. At a two-torsion centre put instead
\(y=\rho U\), \(x=a+\rho^2W\), obtaining
\[
 U^2=(3a^2-1)W+3a\rho^2W^2+\rho^4W^3.
                                                        \tag{A.62.16}
\]
The first nonzero \(x-a\) degree is \(k/2\).
The same valuation estimates with \(j\) replaced
by \(2j\) leave a unit multiple of \(U^k\).

In both cases the special equations have unit
partials in \(W,V\), and the resulting curve is
\(V^p-V=cU^k\), \(c\ne0\). These are integral
smooth algebraic charts whose generic coordinate
changes are invertible on the actual covering field.
The special curve is geometrically integral since
\(k\) is prime to \(p\). Over the algebraic residue
field its nonzero coefficient can be absorbed in \(U\).
At infinity, with \(t=U^{-1}\), one has
\(v(t)=p\) and \(v(V)=-k\). As \(k\mid p-1\),
the element \(tV^{(p-1)/k}\) is a uniformizer;
its difference under a nonzero translation of \(V\)
has valuation \(k+1\). The different is
\((p-1)(k+1)\), and Hurwitz gives genus
\((p-1)(k-1)/2\). The one-branch-point
Deuring--Shafarevich formula gives special \(p\)-rank
zero for every tail \cite[Corollary 1.8]{fghCrew1984}.

\paragraph{Genus exhaustion, the actual map, and all marked points.}
The characteristic-zero division divisor, translation
lifts and quarter-rotation argument of A.61 identify
the actual deck group as the exponent-\(p\) Heisenberg
group and give the asserted degree, cusps and genus.
No feature of that argument requires inertness.

The positive-genus main and all tail valuations just
constructed persist in the stable model, by the same
common-regular-resolution argument as A.58 and A.61,
using \cite[Theorem 9.2.2, Remark 9.2.6 and
Theorem 9.3.21]{fghLiu}. They are distinct since their
restrictions to the elliptic model have distinct
finite centres, or its generic point for the main.
Their genera sum to
\[
 1+p\frac{3(p-1)}2+p(p-3)\frac{p-1}2
                  =1+\frac{p^2(p-1)}2.                \tag{A.62.17}
\]
The semistable genus formula therefore forces compact
type and makes every remaining component rational.
The actual map \(q\) extends to the good elliptic
model after the same regular resolution and descent
of rational exceptional curves. The main component
already maps with degree \(p\), so the degree of a
pulled-back ample line bundle makes every other
component map constantly. Each off-main subtree
has one attaching point and the same constant image.

The ordinary CM factorization splits the generic
\(p\)-torsion into the etale factor of order \(p\)
and the connected factor of order \(p\). The first
specializes bijectively to \(H\), and the second
specializes to \(O\). Thus exactly \(p\) distinct
generic points of \(E[p]\) specialize to each point
of \(H\). Each has one ramified cusp on \(C_p\).
The main radicial map has one point above each
point of \(H\), while all positive-genus tails
and their connecting subtrees lie over centres
outside \(H\). Separating the \(p\) marks at
each main point therefore puts them on a rational
marked tree. No mark stays at its attachment node.
Every selected cusp is joined to another by a path
consisting entirely of rational components, including
both endpoints.

Statement A.42 now gives full fixed inertia at every
coefficient prime and permitted packet, and A.59
gives all normalizers. The accompanying exact
certificate checks the division differential,
the actual separable-isogeny descent and every
critical multiplicity at several ordinary primes.
The unramified translation comparison
(A.62.5)--(A.62.10) proves the assertion uniformly.
\end{proof}
